\chapter{Analýza a~specifikace požadavků}
\label{ch:analyse}
Analýza a~specifikace požadavků tvoří most mezi reálným prostředím a~technickou realizací informačního systému.
Cílem této kapitoly je porozumět každodennímu provozu ubytovacího zařízení\,--\,od rezervačních procesů přes správu jednotlivých ubytovacích jednotek až po doplňkové služby\,--\,a~na tomto základě formulovat požadavky na nový systém.
Důraz je kladen na přesnou identifikaci uživatelských rolí, neboť ty vymezují kompetence i~interakce v~rámci budoucí aplikace a~zajišťují, že výsledné řešení bude odpovídat nárokům personálu i~očekávání hostů.

Závěrečná část kapitoly se věnuje návrhu doménového modelu, který prostřednictvím základní datové struktury vizualizuje klíčové entity a~jejich vzájemné vazby a~tvoří základ pro pozdější datový model.

\section{Současná situace Kempu Olšovec}
Kemp Olšovec je rekreační areál ležící na území městyse Jedovnice na hranici Chráněné krajinné oblasti Moravský kras.
Hostům nabízí širokou paletu ubytovacích možností\,--\,přijet je možné se stanem, osobním automobilem, obytným vozem či karavanem, případně lze využít některou z~ubytovacích jednotek přímo v~areálu.
Ten zahrnuje turistické chaty v~původním i~modernizovaném provedení, dřevěné chaty s~vyšším standardem výbavy a~kompletně vybavené zděné chaty a~bungalovy.

Areál zaujímá rozlohu přibližně 8{,}5~hektaru a~rozprostírá se podél břehu rybníka Olšovec, jehož vodní plocha o~výměře 40~hektarů je jedním z~hlavních důvodů návštěvnosti lokality.
Celková ubytovací kapacita dosahuje až 1\,500~míst v~různých formách ubytování, přičemž pevná lůžková kapacita v~chatách, bungalovech a~pokojích zděné budovy činí přibližně 370~lůžek.
V~hlavní sezóně (červenec a~srpen) je obsazenost areálu prakticky stoprocentní\,--\,počet pokojodnů\footnote{Pokojoden vyjadřuje obsazení jedné ubytovací jednotky po dobu jednoho dne.} v~tomto období dosahuje téměř 4\,000.
Celkový roční počet poskytnutých noclehů se pohybuje kolem 45\,000 při přibližně 21\,000 ubytovaných hostech\footnote{Provozní data byla poskytnuta vedením kempu Olšovec a~reflektují stav v~roce 2025.}.
Bungalovy a~pokoje ve zděné budově jsou v~provozu celoročně, zatímco chatky, stanové plochy a~místa pro karavany jsou zpřístupněny od dubna do října.
Tato sezónnost přímo ovlivňuje požadavky na systém, který musí zvládat výrazné výkyvy zátěže\,--\,od minimálního provozu v~zimních měsících po plnou obsazenost v~letní sezóně.

Cenová politika kempu je nastavena s~ohledem na široké spektrum návštěvníků a~řadí areál mezi cenově dostupná ubytovací zařízení v~Jihomoravském kraji.
Atraktivní volbou je zejména pro rodiny s~dětmi, mladé páry a~skupinové pobyty.
Provozovatel zveřejňuje samostatné ceníky pro jednotlivce, skupiny a~školy, což se musí promítnout do struktury ceníků v~novém systému.
Areál disponuje vlastní jídelnou se zázemím pro plnou penzi i~polopenzi, společenskými prostory a~sportovním vybavením, díky čemuž zde lze realizovat školní výlety, vzdělávací kurzy, adaptační pobyty, sportovní soustředění i~firemní akce pro početné skupiny.
Vedle vlastního ubytovacího provozu se v~kempu i~v~jeho bezprostředním okolí konají hudební a~kulturní akce, závody horských kol na navazujících tratích Singletrail Moravský kras a~další pravidelné události, které dále posilují návštěvnost lokality.

Samostatnou skupinu klientely tvoří zahraniční hosté, mezi nimiž tradičně převažují návštěvníci z~Německa, Nizozemska, Slovenska, Belgie a~Polska.
Jejich podíl na celkovém počtu ubytovaných v~hlavní sezóně dosahuje přibližně 14~\% všech hostů a~75~\% všech cizinců.
Hlavním důvodem návštěvnosti této skupiny je strategická poloha kempu na okraji Moravského krasu, který svými jeskyněmi (Balcarka, Kateřinská, Punkevní, Sloupsko-šošůvská), propastí Macocha a~hustou sítí cyklotras i~turistických tras patří k~nejnavštěvovanějším přírodním lokalitám v~České republice~\cite{ochranaprirody-moravskykras}.
Hosté zde nacházejí pohodlnou základnu pro výlety do Brna (25~km), Blanska (11~km) i~k~okolním kulturním památkám a~chráněným územím.
Přímo v~areálu a~jeho nejbližším okolí je k~dispozici koupání v~rybníce, trasy Singletrail v~přilehlém lese, půjčovna šlapadel a~paddleboardů, vodní lyžování, minigolf, plážový volejbal, dětské hřiště, trasy pro horská kola a~v~letních měsících také rybolov.
Olšovec tak nepředstavuje pouze přechodné ubytovací zařízení, ale plnohodnotné rekreační středisko sezónního cha\-ra\-kte\-ru, což klade na rezervační i~provozní systém vyšší nároky než u~tradičního ubytovacího zařízení.

\subsection{Práce v~kempu z~pohledu recepčních}

Tato sekce představuje pracovní náplň recepčního, který v~rámci ubytovacího procesu plní centrální koordinační roli a~zodpovídá za plynulý chod areálu.
Jeho činnost lze rozdělit do několika tematických okruhů, jež zhruba odpovídají typickému průběhu směny\,--\,od odbavení přijíždějících hostů přes správu vjezdových karet a~check-out až po závěrečnou denní uzávěrku.
Souběžně s~těmito úkony probíhá vyřizování rezervací a~průběžná komunikace s~hosty.

Primárním úkonem je odbavení hosta, jež zahrnuje evidenci, platbu a~předání informací o~pobytu.
V~první fázi se zjišťuje typ a~forma požadovaného ubytování a~existence případné rezervace.
Od všech ubytovaných osob jsou poté vyžadovány doklady totožnosti, na jejichž základě probíhá evidence v~ubytovacím systému.
Do něj se zadávají další údaje\,--\,registrační značka vozidla, konkrétní typ ubytovací jednotky (zděná chata, dřevěná chata, bungalov, karavan, obytný vůz, stan apod.) a~doložky promítající se do finální ceny, jako jsou rekreační poplatky a~poplatky za osoby a~vozidla, jejichž výše se liší podle typu ubytování.
Po úhradě se tiskne doklad, jehož kopii host podepisuje.
U~cizinců přibývá podpis evidenčního formuláře pro cizineckou policii.
Závěrem odbavení se předávají klíče a~další informace o~ubytování\,--\,provozní doba, postup při opouštění areálu, vybavení chaty i~režim sprch.

Na odbavení navazuje systém vjezdových karet do areálu.
Host si může u~recepce, obvykle při příjezdu, vyzvednout kartu, prostřednictvím které během pobytu samostatně reguluje vjezd a~výjezd.
Recepční vybírá vratnou kauci a~v~samostatném systému kartu aktivuje na sjednanou dobu pobytu, přičemž zaznamenává i~výši kauce.

Ranní směna obstarává check-out odjíždějících hostů.
Prvním krokem je vyplnění tzv.\ úklidového formuláře, který na základě ubytovacího plánu identifikuje chaty určené k~úklidu daného dne a~od nichž je nutné do desáté hodiny dopolední převzít klíče.
Vytištěný formulář se předává úklidové směně.
V~případě, že host povinnost vrácení klíčů zakotvenou v~provozním řádu kempu nedodrží, řeší recepční situaci operativně tak, aby mohl úklid v~plánovaném čase proběhnout.
Do agendy ranní směny dále spadá řešení reklamací, stížností a~případných škod na ubytovacích jednotkách či jiném majetku v~areálu.

Stěžejním úkonem konce směny je denní uzávěrka, tedy souhrnné vyrovnání tržeb za daný den.
Vytváří se ve dvou kopiích: jedna putuje v~obálce do kanceláře vedení kempu, druhá zůstává na recepci spolu s~účtenkami za hotovostní a~bezhotovostní platby a~s~doklady o~ubytovaných cizincích.
Do téže obálky se přikládá také denní evidence tržeb vytištěná ze stejného systému.
Ve stejném systému se poté zaeviduje seznam cizinců ubytovaných během daného dne a~tato evidence se nahraje na webový portál Ubyport, odkud ji načítá cizinecká policie; portál zároveň generuje tzv.~doručenku, kterou recepční e-mailem zasílá vedení kempu.
Veškeré dokumenty spojené s~uzávěrkou včetně evidence tržeb jsou archivovány v~elektronické podobě (PDF) v~recepčním počítači a~jejich kopie se z~recepčního e-mailu odesílá hlavní účetní.

Souběžně s~výše popsanými úkony probíhá během směny také vyřizování rezervací.
Týkají se nejčastěji všech typů chat, prostor pro karavany a~obytné vozy, firemních setkání a~prostor pro oslavy.
Komunikace s~hosty probíhá zpravidla telefonicky či e-mailem.
Rezervace se ukládá do ubytovacího portálu s~uvedením jména a~kontaktu hosta a~v~případě online objednávky též údajů z~rezervačního formuláře\footnote{viz \url{https://www.olsovec.cz/rezervace/}}.
Mezi další průběžné povinnosti patří řešení storen, předávání požadavků pracovníkům údržby, dohled nad dodržováním provozního řádu a~poskytování informací o~areálu i~službách v~jeho okolí.

\section{Neformální specifikace požadavků}
\label{sec:pozadavky}
Turistický kemp Olšovec,~s.r.o. se z~důvodu provozní neefektivity rozhodl implementovat nový centrální informační systém.
Současná správa rezervací pomocí roztříštěných nástrojů (aplikace Excel ze sady nástrojů MS Office, e-mail, telefon, osobní domluvy s~hosty přímo na recepci, samostatný software používaný k~vystavování dokladů a~evidenci osob) vede k~vysoké chybovosti, chybějícím nebo duplicitním rezervacím a~pomalému odbavení.
Cílem je zavedení jednotného systému do půl roku, přičemž minimální funkcionalita v~podobě vyřizování rezervací musí být k~dispozici od nového roku z~důvodu plynulého přechodu na nový systém a~seznamování zaměstnanců se systémem během nevytíženého zimního období.

Stěžejním prvkem systému je veřejně dostupný online \textbf{rezervační formulář} nacházející se na webových stránkách kempu.
Každá rezervovatelná jednotka (chata, karavanové stání, stanové místo) zde obsahuje základní parametry (kapacita, přípojka elektřiny, bezbariérový přístup,\,\ldots) a~odkaz na prezentaci, která se nachází na webových stránkách.
Kdokoliv zde může provést rezervaci ve zvolený termín na jednu rezervovatelnou jednotku.
Každou takto vytvořenou rezervaci musí recepční \textbf{potvrdit}.

Systém ovšem musí umět pracovat i~se \textbf{složitějšími rezervacemi} (např. rezervace více míst, skupinové rezervace, stěhování mezi ubytovacími jednotkami v~průběhu pobytu).
Tyto specifické rezervace není potřeba řešit v~rezervačním formuláři pro hosty, neboť tvoří přibližně 2~\% z~celkového množství rezervací a~přistupuje se k~nim individuálně.
Jejich vyřízení je proto ponecháno na recepčních.
U~větších rezervací musí systém navíc umožnit \textbf{objednání stravy}, kterou pak recepce předá kuchyni.

Z~důvodu zákonné povinnosti evidovat osobní údaje všech ubytovaných hostů musí systém u~hromadných a~skupinových rezervací umožnit organizátorovi provést \textbf{online check-in} všech ubytovaných hostů, případně zadat registrační značky vozidel, která se budou v~kempu vyskytovat.
Příjezd například 150 osob v~krátkém časovém úseku by jinak recepci zcela ochromil.
Účastníci školních výletů u~sebe často doklad totožnosti nemají nebo jím vůbec nedisponují, a~tato funkcionalita tak zároveň řeší praktický provozní problém.
Osobní údaje hostů musí být přehledně dohledatelné v~digitální \textbf{knize ubytovaných hostů} v~souladu s~platnou legislativou.

V~případě dražšího pobytu vystavují recepční podklady pro účetní, aby v externím účetním softwaru vystavila \textbf{fakturu}, kterou host hradí část ceny z pobytu.
Fakturu účetní zasílá hostovi na e-mail a její stav průběžně eviduje v systému.
Výsledný zúčtovací doklad, vystavený recepčním v~den příjezdu po kompletním odbavení hosta na recepci, je automaticky ponížen o cenu položek již uhrazených prostřednictvím faktury.

Omezený přístup do systému náleží \textbf{provoznímu personálu} (uklízečky, údržbáři), který bude využívat \textbf{zjednodušenou mobilní aplikaci}.
Uklízečky zde uvidí denní plán úklidů s~informací, ke kterým chatám již byly vráceny klíče.
Po dokončení úklidu mění stav jednotky tak, aby bylo recepčním zřejmé, že je již možné vydat klíče dalšímu hostovi.
Údržbáři v~aplikaci vidí nahlášené závady, které do systému zadávají recepční na základě hlášení od hostů.

Přístup do systému musí mít také \textbf{účetní} společnosti, jež veškeré vystavené doklady zaúčtovává.
Systém jí musí umožnit nahlížet na všechny \textbf{vydané doklady} (účtenky, faktury), \textbf{přijaté platby} a~\textbf{peněžní deník}.
Nezbytné jsou rovněž \textbf{denní uzávěrky} a~další finanční sestavy, které systém musí poskytovat.
Všechny tyto doklady musí být v~souladu s~platnou legislativou, přehledné a~strukturované.

Se systémem dále pracuje \textbf{vedoucí kempu}, který sleduje přehled obsazenosti, nahlíží na účetní přehledy a~vyhodnocuje provozní a~ekonomické statistiky areálu.
Spravuje rovněž \textbf{ceníky}, \textbf{položky možné k~zaúčtování} a~\textbf{rezervovatelné jednotky} (jejich dostupnost, kapacitu a~parametry).
Do jeho kompetencí spadá také \textbf{správa zaměstnanců} včetně přidělování jejich systémových rolí.

Nejintenzivněji budou se systémem pracovat \textbf{recepční}.
Vykonávají klasickou recepční práci a~navíc \textbf{kontrolu hostů a~vozidel}, kteří se nenacházejí v~chatách (stany, osobní automobily, karavany, obytné vozy,\,\ldots).
Pro tuto kontrolní činnost budou používat mobilní aplikaci umožňující načítání QR~kódů z~nalepovacích štítků (hosté tyto štítky obdrží při check-inu na recepci) i~skenování registračních značek vozidel, případně jejich ruční zadání.

Při \textbf{klasické recepční práci} pracují recepční pouze s~desktopovým rozhraním, kde především přijímají rezervace, provádějí odbavení hostů a~obsluhují procesy popsané výše.
Pro recepční je rovněž důležitá \textbf{integrace s~externími systémy}: systémem pro prokazování totožnosti \texttt{E-doklady}, viz \ref{sec:edoklady}, systémem pro ohlašování pobytu cizinců \texttt{Ubyport}, viz \ref{sec:ubyport}, a dalšími.
Některé z~těchto systémů aktuálně nejsou v~plném provozu, případně se očekává jejich legislativní úprava; architektura nového systému proto musí být navržena tak, aby umožňovala jejich snadné napojení či úpravu napojení do budoucna.

Z~hlediska \textbf{nefunkčních požadavků} jsou na systém kladeny následující nároky.
\textbf{Uživatelské rozhraní} musí být přehledné a~moderní; klíčové ovládací prvky\,--\,zejména \uv{plachta} obsazenosti, tedy tabulkové znázornění rezervací jednotek v~čase\,--\,však musí dodržovat konvence známé z~obdobných ubytovacích systémů, aby si recepční rychle osvojili práci s~novým systémem a~orientace v~něm zůstala intuitivní.
Vzhledem k~významnému podílu zahraničních hostů uvedenému v~úvodu kapitoly musí být zajištěna \textbf{lokalizace} uživatelského rozhraní rezervačního formuláře i~vystavovaných dokladů.
Systém musí zvládat \textbf{souběžnou práci} všech recepčních a~pracovníků provozního personálu v~hlavní sezóně, a~současně musí být schopen obsloužit hosty.
Vystavované doklady, evidence ubytovaných hostů i~zpracování osobních údajů musí splňovat požadavky \textbf{platné legislativy}.
Implementace systému musí v~maximální možné míře využít \textbf{stávající infrastrukturu} kempu (recepční počítače, tiskárny, síťové prvky, čtečky), čímž se minimalizují pořizovací náklady, a~zároveň musí být zajištěna snadná přenositelnost serverové části a~nasazení na další pracovní stanice.
Architektura systému musí být \textbf{rozšiřitelná}, a~to především s~ohledem na legislativní nestabilitu, napojení na externí státní systémy; změny v~těchto integracích musí být proveditelné bez zásadních zásahů do zbytku systému.

\section{Uživatelské role}
\label{sec:role}
Se systémem budou pracovat uživatelé v~následujících rolích:
\begin{itemize}
    \item \textbf{Host}\,--\,osoba, která má zájem o~vytvoření rezervace, případně organizátor skupinové rezervace; využívá veřejný rezervační formulář a~provádí online check-in členů své skupiny; nemá vytvořený uživatelský účet v~systému.
    \item \textbf{Recepční}\,--\,hlavní denní operátor systému; v~desktopovém rozhraní spravuje rezervace, manuálně zadává složité a~skupinové rezervace, provádí check-in a~check-out hostů, eviduje ubytované a~vystavuje doklady; rovněž zadává požadavky na údržbu a~úklid a~prostřednictvím mobilní aplikace provádí kontrolu hostů a~vozidel v~areálu (skenování registračních značek a~QR~kódů ze štítků).
    \item \textbf{Účetní}\,--\,osoba zodpovědná za zaúčtování vystavených dokladů; přistupuje k~vydaným dokladům, přijatým platbám a~peněžnímu deníku, pracuje s~denními uzávěrkami a~dalšími finančními sestavami a~exportuje data do účetního softwaru.
    \item \textbf{Uklízečka}\,--\,provozní pracovník využívající mobilní aplikaci, ve které vidí denní plán úklidů (včetně informace o~vrácení klíčů) a~po dokončení práce zaznamenává úklid každé jednotky.
    \item \textbf{Údržbář}\,--\,provozní pracovník využívající mobilní aplikaci; přijímá v~ní požadavky na opravy, které do systému zadali recepční na základě hlášení od hostů, a~po vyřešení požadavek v~aplikaci uzavírá.
    \item \textbf{Vedoucí kempu}\,--\,sleduje celkovou obsazenost, finanční přehledy a~provozní statistiky; spravuje konfigurační nastavení systému (ceníky, položky k~zaúčtování, dostupnost a~parametry rezervovatelných jednotek); obdobně jako recepční spravuje rezervace; zajišťuje správu zaměstnanců a~přidělování jejich systémových rolí.
\end{itemize}

\subsection{Činnosti jednotlivých uživatelů se systémem}
Jednotlivé uživatelské role v~systému provádějí následující činnosti:
\begin{itemize}
    \begin{samepage}
        \item \textbf{Host}
        \begin{itemize}
            \item prohlížení volných kapacit v~rezervačním formuláři,
            \item vyplnění a~odeslání jednoduché rezervace přes webový formulář,
            \item zobrazení informací o~své rezervaci,
            \item provedení online check-inu (zadání osobních údajů a~registračních značek vozidel),
            \item rozeslání odkazu na online check-in členům skupiny (v~případě skupinové rezervace),
            \item zrušení rezervace.
        \end{itemize}
    \end{samepage}

    \item \textbf{Recepční}
    \begin{itemize}
        \begin{samepage}
            \item \underline{kontrola areálu}
            \begin{itemize}
                \item skenování QR~kódů z~nalepovacích štítků,
                \item načítání registračních značek zaparkovaných vozidel,
                \item manuální vyhledání hostů nebo vozidel.
            \end{itemize}
        \end{samepage}

        \begin{samepage}
            \item \underline{zajištění provozu}
            \begin{itemize}
                \item vytváření rozpisů úklidu,
                \item zadávání nahlášených závad,
                \item úprava dostupnosti chat (např. z~důvodu poruch),
                \item plánování objednávek a výdejů stravy.
            \end{itemize}
        \end{samepage}

        \begin{samepage}
            \item \underline{správa rezervací}
            \begin{itemize}
                \item potvrzování webových rezervací,
                \item vytváření složitějších rezervací,
                \item úprava a~rušení rezervací,
                \item objednávání stravy pro skupinové rezervace,
                \item komunikace s~hosty.
            \end{itemize}
        \end{samepage}

        \begin{samepage}
            \item \underline{ubytování hostů a~finanční operace}
            \begin{itemize}
                \item evidence ubytovaných hostů a~registračních značek vozidel v~areálu,
                \item vystavování podkladů pro účetní,
                \item vystavování účtenek,
                \item tisk kontrolních nalepovacích štítků,
                \item vystavování opravných daňových dokladů.
            \end{itemize}
        \end{samepage}

        \begin{samepage}
            \item \underline{administrativní operace}
            \begin{itemize}
                \item odesílání zákonných údajů do státních systémů dle platné legislativy,
                \item vytváření denních uzávěrek.
            \end{itemize}
        \end{samepage}
    \end{itemize}

    \begin{samepage}
        \item \textbf{Účetní}
        \begin{itemize}
            \item kontrola vydaných dokladů,
            \item evidence stavu úhrad vystavených faktur a~jeho předávání recepčním,
            \item stahování denních účetních závěrek,
            \item stahování finančních sestav a~přehledů,
            \item export dat do účetního softwaru.
        \end{itemize}
    \end{samepage}

    \begin{samepage}
        \item \textbf{Uklízečka}
        \begin{itemize}
            \item zobrazení denního plánu úklidu,
            \item náhled na stav jednotlivých chat (včetně informace o~vrácení klíčů),
            \item změna stavu chaty po dokončení úklidu.
        \end{itemize}
    \end{samepage}

    \begin{samepage}
        \item \textbf{Údržbář}
        \begin{itemize}
            \item zobrazení seznamu nahlášených závad,
            \item náhled na detail jednotlivých požadavků,
            \item uzavření požadavku po vyřešení závady.
        \end{itemize}
    \end{samepage}

    \item \textbf{Vedoucí kempu}
    \begin{itemize}
        \item správa rezervací ve stejném rozsahu jako recepční (viz výše),
        \item nahlížení na strategické statistiky a~finanční přehledy,
        \begin{samepage}
            \item \underline{správa systému}
            \begin{itemize}
                \item správa uživatelských účtů a~přiřazování rolí,
                \item správa ceníků a~položek k~zaúčtování,
                \item konfigurace rezervovatelných jednotek (dostupnost, kapacita, parametry),
                \item konfigurace dalších systémových nastavení.
            \end{itemize}
        \end{samepage}
    \end{itemize}
\end{itemize}

\subsection{Diagram případů užití}
Diagram případů užití, který vychází z~rolí a~činností uvedených v~předchozích 2 sekcích, je znázorněn na obrázcích \ref{fig:uc-guest} a~\ref{fig:uc-receptionist}.

\begin{figure}[H]
    \centering
    \includegraphics[width=1.0\textwidth]{uc-1.pdf}
    \caption{Diagram případů užití\,--\,hosté a~správa kempu.}
    \label{fig:uc-guest}
\end{figure}

\begin{figure}[H]
    \centering
    \includegraphics[width=1.0\textwidth]{uc-2.pdf}
    \caption{Diagram případů užití\,--\,recepční, vedoucí.}
    \label{fig:uc-receptionist}
\end{figure}

\section{Doménový model}
Doménový model, znázorněný na obrázku~\ref{fig:domain-model}, představuje stěžejní artefakt fáze analýzy požadavků.
Tvoří konceptuální vrstvu mezi neformální specifikací zadavatele a~technickým návrhem systému prostřednictvím formalizace problémové domény do podoby entit a~jejich vzájemných vztahů.

Zařazení doménového modelu do analytické části práce je metodologicky odůvodněné.
Model vznikal iterativní konzultací s~cílovými uživateli systému, zejména s~vedením kempu a~zkušenými recepčními; po každém návrhu byly identifikovány nepřesnosti a~model byl upravován tak, aby věrně odpovídal reálným provozním procesům kempu, a~nikoliv pouze jejich vývojářské interpretaci.

Pro vytvoření doménového modelu v~této fázi vývoje hovoří několik důvodů.
Slouží jako sdílený konceptuální jazyk mezi vývojářem a~zadavatelem, neboť terminologie zachycená v~entitách a~vazbách se promítá do uživatelského rozhraní, dokumentace i~budoucí komunikace nad systémem, a~obě strany ji proto musí interpretovat shodně.
Vizuální podoba modelu dále umožňuje netechnickým účastníkům konzultace ověřit věrnost zachycení provozu kempu, což u~rozsáhlejších textových specifikací zpravidla není snadno dosažitelné.
Modelování rovněž odhaluje rozpory či neúplnosti v~požadavcích, jejichž náprava je řádově méně nákladná než v~průběhu návrhu či implementace, a~současně tvoří východisko pro pozdější návrh datové vrstvy.
Pro lepší orientaci jsou entity rozděleny do čtyř barevně odlišených skupin odpovídajících hlavním funkčním okruhům systému.

Oranžově vyznačená skupina tvoří jádro rezervační logiky a~evidence ubytovaných hostů.
Mapuje vztahy mezi tím, kdo ubytování poptává (entity \texttt{Reservation} a~\texttt{GroupReservation}), kdo je skutečně ubytován (\texttt{Guest}), na jakém místě (\texttt{Spot}, \texttt{SpotGroup}) a~kdo za pobyt platí.
Skupinové rezervace umožňují \uv{předrezervaci} více míst najednou s~nutností dodatečného ověření před vytvořením rezervace již předrezervovaných jednotek, čímž je pokryt scénář komplexních rezervací popsaný v~sekci~\ref{sec:pozadavky}.
Entita \texttt{Vehicle} reflektuje praktickou potřebu evidovat registrační značky pro kontrolu vjezdu a~parkování, pomocí entity \texttt{Nationality} je pak evidována státní příslušnost ubytovaných hostů.
Entita \texttt{ReservationServiceItem} propojuje rezervaci s~jednotlivými službami, které se k~ní vážou (např. objednanou stravou pro skupinové rezervace či doplňkovými poplatky).
Entita \texttt{AccessCard} je použita pro ukládání informací o fyzických kartách na ovládání závor v~areálu a výše přijaté kauce.

Žlutozelená skupina reprezentuje operativní a~personální vrstvu systému, jež se uplatňuje především při práci zaměstnanců v~terénu a~při správě fyzických prostředků areálu.
Entity \texttt{CleaningPlan} a~\texttt{CleanInfo} zajišťují plynulý cyklus přípravy ubytovacích kapacit\,--\,od plánování úklidu přes zaznamenání jeho dokončení až po informaci, že jednotka může být předána dalšímu hostovi.
Entita \texttt{MaintenanceIssue} eviduje nahlášené závady; vazba na entitu \texttt{User} (vyhrazená pro pracovníky provozního personálu) zachycuje, který konkrétní zaměstnanec danou závadu vyřešil, čímž je zajištěna dohledatelnost údržbářských zásahů.
Entita \texttt{OutOfOrder} slouží k~označení dočasně nedostupných jednotek (např. z~důvodu poruchy nebo plánované údržby), a~brání tak chybnému ubytování hostů na neprovozuschopných místech.
Entita \texttt{Event} reprezentuje akce konané v~areálu a~slouží k~informování potenciálních hostů již při vytváření rezervace, čímž se předchází stížnostem hostů, kteří by jinak nemuseli být na konání akcí v~areálu připraveni.

Růžová skupina pokrývá ekonomickou agendu a~definici nabízených produktů kempu.
Tato vrstva transformuje poskytnuté služby do finanční podoby.
Entita \texttt{Service} představuje jednotku poskytované služby, pro niž jsou dále specifikovány \texttt{ServiceType} (kategorie), \texttt{ServiceText} (textový popis na dokladech) a~\texttt{VatRate} (sazba DPH).
Samostatnou entitou je \texttt{Meal}, která zachycuje stravovací položky vázané na rezervaci s fixními typy jídel.
Model rozlišuje mezi účtenkami (\texttt{Bill}) a~fakturami (\texttt{Invoice}); účtenky slouží jako zúčtovací doklady, faktury jsou určeny především pro předčasné platby\,--\,\uv{zálohy} na pobyt.
Účtovatelné položky na účtenkách jsou prostřednictvím entity \texttt{BillItem} navázány přímo na entitu \texttt{Service}, čímž je zajištěna konzistence mezi definovanou nabídkou kempu a~jejím promítnutím do vystavovaných dokladů.
Sebereferenční vazba u~účtenky umožňuje realizaci opravných daňových dokladů.
Entita \texttt{Language} umožňuje vystavování dokladů v~jazyce hosta, což je relevantní zejména pro významný podíl zahraničních klientů zmíněný v~úvodní sekci kapitoly.
\texttt{FinancialClosing} reprezentuje účetní závěrky, díky nimž je zajištěna transparentnost finančních toků a~předání podkladů hlavní účetní.

Šedé entity tvoří průřezovou vrstvu autentizace a~oprávnění uživatelů systému.
Entita \texttt{User} reprezentuje konkrétního uživatele pracujícího se systémem, entita \texttt{Role} definuje jeho systémovou roli odpovídající rolím popsaným v~sekci~\ref{sec:role}.
Entita \texttt{Passkey} obsahuje autentizační informace uživatele pro přihlášení do systému.

\begin{figure}[H]
    \centering
    \includegraphics[width=0.95\textwidth]{domain-model.pdf}
    \caption{Doménový model informačního systému.}
    \label{fig:domain-model}
\end{figure}
